\section{TestRail Reflection}

\subsection{Testing Overview}

During Snapshot 4, the QA team focused on final system validation, regression testing, accessibility reviews, performance optimization, and production release preparation for the AirWatch platform.

Testing covered all major project components developed throughout previous snapshots, including:
\begin{itemize}
    \item AQI monitoring
    \item Interactive maps
    \item Saved locations
    \item Push notifications
    \item Health advisory recommendations
    \item Cross-platform mobile and web compatibility
\end{itemize}

The final QA phase emphasized release stability and user experience consistency before deployment.

\subsection{Test Execution Summary}

\begin{table}[h]
\centering
\begin{tabular}{|l|c|}
\hline
\textbf{Metric} & \textbf{Result} \\
\hline
Total Test Cases Written & 92 \\
Executed & 92 \\
Passed & 86 \\
Failed & 4 \\
Blocked & 2 \\
Overall Pass Rate & 93\% \\
\hline
\end{tabular}
\end{table}

\subsection{Major Issues Identified}

\subsubsection{Dark Mode UI Inconsistencies}

Several mobile screens displayed incorrect text contrast after enabling dark mode.

Examples included:
\begin{itemize}
    \item unreadable advisory feed text
    \item overlapping buttons on smaller Android devices
    \item inconsistent icon coloring
\end{itemize}

The issue was partially resolved through updated theme styling and additional responsive layout testing.

\subsubsection{Performance Issues During Heavy Map Usage}

Rapid zooming and panning on the AQI heat-map occasionally caused lag on older devices.

The development team improved performance by:
\begin{itemize}
    \item reducing overlay refresh frequency
    \item optimizing map rendering
    \item limiting unnecessary API requests
\end{itemize}

\subsubsection{Accessibility Audit Findings}

Accessibility testing identified several areas needing improvement:
\begin{itemize}
    \item missing screen reader labels
    \item insufficient contrast ratios
    \item incomplete keyboard navigation support on web builds
\end{itemize}

Most accessibility issues were resolved before final QA approval.

\subsection{QA Reflection}

Using TestRail throughout development significantly improved testing organization and regression tracking. Maintaining separate test suites for maps, notifications, and advisory systems helped isolate failures quickly and reduced duplicate testing effort.

The team learned that:
\begin{itemize}
    \item cross-platform testing should begin earlier in development
    \item accessibility validation should occur continuously instead of only during final QA
    \item performance testing on low-end devices is critical for mobile stability
\end{itemize}

The addition of automated regression testing also reduced the number of repeated manual verification steps during final release preparation.

\subsection{Process Improvements}

For future releases, the QA team recommends:
\begin{itemize}
    \item earlier integration of automated UI testing
    \item expanded accessibility testing coverage
    \item additional stress testing for notification services
    \item more physical device testing across Android versions
    \item automated dark mode validation
\end{itemize}

The team also plans to improve CI/CD workflows using deployment automation and crash monitoring tools.

\subsection{Final QA Assessment}

The Snapshot 4 build successfully met release readiness requirements. Core features operated reliably across supported platforms, and most critical defects were resolved before deployment.

The QA team approved the application for production release following completion of regression and accessibility verification.
